\section{Hypergraph manipulations}

\subsection{Bethe hypergraphs}

\begin{definition}[Delta tensor on generic semirings]
    Given a semiring $\gensemiring=\gensemiringexp$ and a collection of variables $\catvariableof{\edge}$ of same dimension, the corresponding delta tensor $\identityat{\catvariableof{\edge}}$
    \begin{align*}
        \identityat{\indexedcatvariableof{\edge}}
        = \begin{cases}
              \srone & \ifspace \uniquantwrtof{\node,\secnode\in\edge}{\catindexof{\node} = \catindexof{\secnode}} \\
              \srzero & \elsetext
        \end{cases} \, .
    \end{align*}
\end{definition}

The main manipulation operation with delta tensor is shown in the next lemma.

\begin{lemma}
    Given a tensor network $\extnetwith$ with respect to a semiring $\gensemiring=\gensemiringexp$.
    For any set of edges $\secedges\subset\edges$ each containing a variable $\node$ define cores $\sechypercoreofat{\edge}{\catvariableof{\edge\backslash\{\node\}},\catvariableof{\nodeof{\edge}}} $ with coordinates
    \begin{align*}
        \sechypercoreofat{\edge}{\indexedcatvariableof{\edge\backslash\{\node\}},\catvariableof{\nodeof{\edge}}=\catindexof{\node}}
        = \hypercoreofat{\edge}{\indexedcatvariableof{\edge}} \, .
    \end{align*}
    Then we have
    \begin{align*}
        \extnetwith =
        \srcontractionof{
            \{\hypercoreofat{\edge}{\catvariableof{\edge}}\wcols \edge\notin\secedges\}
            \cup
            \{ \sechypercoreofat{\edge}{\catvariableof{\edge\backslash\{\node\}},\catvariableof{\nodeof{\edge}}} \wcols \edge\in\secedges\}
            \cup
            \{\identityat{\catvariableof{\node},\catvariableof{\nodeof{\edge}}\wcols\edge\in\secedges}\}
        }{
            \nodevariables
        } \, .
    \end{align*}
\end{lemma}
\begin{proof}
    Using the neutrality properties of $\srone$ and $\srzero$.
\end{proof}



\begin{definition}
            \glsadd{mp:BetheHypergraph}
        Given a tensor network $\tnetof{\graph}$ on a decorated hypergraph $\graph$, we define the Bethe hypergraph $\secgraph$ as
        $(\secnodes, \secedges \cup \Delta^{\graph})$ where we have
        \begin{itemize}
            \item Recolored Edges $\secedges = \{\tilde{\edge} \wcols \edge\in \edges\}$ where $\tilde{\edge} = \{\node^{\edge} \wcols \node\in\edge\}$, which decoration tensor has same coordinates as $\hypercoreof{\edge}$
            \item Nodes $\secnodes = \nodes \cup (\bigcup_{\edge\in\edges}\tilde{\edge})$ %$\secnodes = \bigcup_{\edge\in\edges}\{\node^{\edge} \wcols \node\in\edge \}$
            \item Delta Edges $\Delta^{\graph} =  \big\{ \{\node\} \cup \{\node^{\edge} \wcols \edge\ni\node \}\wcols \node\in\nodes \big\} $, each of which decorated by a delta tensor $\delta^{\{\node^{\edge} \wcols \edge\ni\node \}}$
        \end{itemize}
\end{definition}

% Bipartite
The resulting overlap graph is bipartite, since any variable appears exactly in one cluster in $\secedges$ and in one cluster of $\Delta^{\graph}$.
This further makes the dual of the Bethe Cluster Hypergraph a proper graph (i.e. edges consistent of node pairs).

We effectively copy each node $\node$ by $\node^{\edge}$.
Let us now show, that when assigning dirac delta tensors the contraction reproduces the original contraction.

By adding delta tensors to each node $\node\in\nodes$ and defining its leg variables by $\node^{\edge}$ for $\edge\in\edges$.
We mark each such delta tensor by a cluster in $\Delta^{\graph}$, as defined in the following (see also \figref{fig:betheDataExample}).

By \lemref{lem:betheInvariance} this construction does not change contractions.

\begin{lemma}[Bethe Invariance]\label{lem:betheInvariance}
    Given a tensor network $\extnetwith$ on a hypergraph $\graph$.
    We define a tensor network on its bethe hypergraph $\secgraph$ by the tensors for $\edge\in\edges$
    \begin{align*}
        \sechypercoreofat{\tilde{\edge}}{\catvariableof{\tilde{\edge}}}
        = \contractionof{\hypercoreofat{\edge}{\catvariableof{\edge}},\identityat{\catvariableof{\edge},\catvariableof{\tilde{\edge}}}}{\catvariableof{\tilde{\edge}}}
    \end{align*}
    and the tensors for $\tilde{\edge}\in\Delta^{\graph}$
    \begin{align*}
        \sechypercoreofat{\tilde{\edge}}{\catvariableof{\tilde{\edge}},\catvariableof{\node}}
        = \identityat{\catvariableof{\tilde{\edge}},\catvariableof{\node}} \, .
    \end{align*}
    Then we have
    \begin{align*}
        \extnetwith =
        \contractionof{\{\sechypercoreofat{\tilde{\edge}}{\catvariableof{\tilde{\edge}}} \wcols \tilde{\edge}\in\secedges \cup \Delta^{\graph}\}}{\nodevariables}
    \end{align*}
\end{lemma}
\begin{proof}
    By invariance of adding dirac deltas.
\end{proof}

Equal constructions:
\begin{itemize}
    \item BP on factor graphs in \cite{mezard_information_2009}
    \item Bethe Cluster graphs in \cite{koller_probabilistic_2009}
\end{itemize}

\begin{figure}[t]
    \begin{center}
        \begin{tikzpicture}[scale=0.35, thick] % , baseline = -3.5pt


    \begin{scope}
        [shift={(23,0)}]

        \node[anchor=center] (text) at (-6,8) {$b)$};

        \draw (2,8) rectangle (4,6);
        \node[anchor=center] (text) at (3,7) {\corelabelsize $\bencodingof{\lor}$};

        \draw (2,5) rectangle (4,3);
        \node[anchor=center] (text) at (3,4) {\corelabelsize $\bencodingof{\land}$};

        \draw (2,2) rectangle (4,0);
        \node[anchor=center] (text) at (3,1) {\corelabelsize $\datacoreof{c}$};

        \draw (2,-1) rectangle (4,-3);
        \node[anchor=center] (text) at (3,-2) {\corelabelsize $\datacoreof{b}$};

        \draw (2,-4) rectangle (4,-6);
        \node[anchor=center] (text) at (3,-5) {\corelabelsize $\datacoreof{a}$};

        \draw (2,-7) rectangle (4,-9);
        \node[anchor=center] (text) at (3,-8) {\corelabelsize $\lambda$};


        \draw[fill] (-3,-8.5) circle (\dotsize);
        \node[anchor=center] (text) at (-4,-8.5) {\colorlabelsize $\indexset$};

        \draw[] (-3,-8.5) to[bend left=-10]  (2,-8);
        \draw[] (-3,-8.5) to[bend left=0]  (2,-5.5);
        \draw[] (-3,-8.5) to[bend left=10]  (2,-2.5);
        \draw[] (-3,-8.5) to[bend left=20]  (2,0.5);


        \draw[fill] (-3,-5.5) circle (\dotsize);
        \node[anchor=center] (text) at (-4,-5.5) {\colorlabelsize $\catvariableof{a}$};

        \draw[] (-3,-5.5) to[bend left=-10]  (2,-4.5);
        \draw[] (-3,-5.5) to[bend left=10]  (2,3.5);

        \draw[fill] (-3,-2.5) circle (\dotsize);
        \node[anchor=center] (text) at (-4,-2.5) {\colorlabelsize $\catvariableof{b}$};

        \draw[] (-3,-2.5) to[bend left=-10]  (2,-1.5);
        \draw[] (-3,-2.5) to[bend left=10]  (2,4);

        \draw[fill] (-3,0.5) circle (\dotsize);
        \node[anchor=center] (text) at (-4,0.5) {\colorlabelsize $\catvariableof{c}$};

        \draw[] (-3,0.5) to[bend left=-10]  (2,1.5);
        \draw[] (-3,0.5) to[bend left=10]  (2,6.5);

        \draw[fill] (-3,3.5) circle (\dotsize);
        \node[anchor=center] (text) at (-4.5,3.5) {\colorlabelsize $\headvariableof{a\land b}$};

        \draw[] (-3,3.5) to[bend left=-10]  (2,4.5);
        \draw[] (-3,3.5) to[bend left=10]  (2,7);

        \draw[fill] (-3,6.5) circle (\dotsize);
        \node[anchor=center] (text) at (-5.5,6.5) {\colorlabelsize $\headvariableof{(a\land b)\lor c}$};

        \draw[] (-3,6.5) to[bend left=10]  (2,7.5);


        \draw[dashed] (-0.5,-12) -- (-0.5,8);

        \node[right] (text) at (0.5,-11) {$\tilde{\edges}$};
        \node[left] (text) at (-1.5,-11) {$\Delta$};

    \end{scope}


    \node[anchor=center] (text) at (-2,8) {$a)$};

    \newcommand{\conposseldec}{3,-5.5}

    \draw[fill] (\conposseldec) circle (\dotsize);
    \draw (\conposseldec) -- (3,-7.5) node[midway, right] {\colorlabelsize ${\indexset}$}; % Unclear, whether this is the best notation!
    \draw[] (2,-7.5) rectangle (4, -9.5);
    \node[anchor=center] (text) at (3,-8.5) {\corelabelsize $\lambda$};

    \draw[-<-] (0,1) -- (0,-1) node[midway,left] {\colorlabelsize $\catvariableof{a}$};
    \draw (-1,-1) rectangle (1, -3);
    \node[anchor=center] (text) at (0,-2) {\corelabelsize $\datacoreof{a}$};
    \draw[-<-] (0,-3) to[bend right=20] (\conposseldec);


    \draw[-<-] (3,1) -- (3,-1) node[midway,left] {\colorlabelsize $\catvariableof{b}$};
    \draw (2,-1) rectangle (4, -3);
    \node[anchor=center] (text) at (3,-2) {\corelabelsize $\datacoreof{b}$};
    \draw[-<-] (3,-3) to[bend right=0]  (\conposseldec);


    \draw[-<-] (6,5) -- (6,-1) node[midway,left] {\colorlabelsize $\catvariableof{c}$};
    \draw (5,-1) rectangle (7, -3);
    \node[anchor=center] (text) at (6,-2) {\corelabelsize $\datacoreof{c}$};
    \draw[-<-] (6,-3) to[bend left=20]  (\conposseldec);


    \draw[-<-] (1.5,5) -- (1.5,3) node[midway,left] {\colorlabelsize $\headvariableof{a \land b}$};
    \draw (-1,3) rectangle (4, 1);
    \node[anchor=center] (text) at (1.5,2) {\corelabelsize $\bencodingof{\land}$};


    \draw[-<-] (3.5,9) -- (3.5,7) node[midway,left] {\colorlabelsize $\headvariableof{(a \land b) \lor c }$};
    \draw (0,7) rectangle (7, 5);
    \node[anchor=center] (text) at (3.5,6) {\corelabelsize $\bencodingof{\lor}$};

\end{tikzpicture}
    \end{center}
    \caption{Example of a Bethe Cluster Graph.
    a) Example of a Tensor Network $\tnetof{\graph}$, which represents the by $\lambda$ averaged evaluation of the formula $(a\land b)\lor c$ on data $\datamap$.
    b) Corresponding Bethe Cluster Hypergraph, which dual is bipartite by the sets $\Delta$ and $\tilde{\edges}$.
    }
    \label{fig:betheDataExample}
\end{figure}

\subsection{Coarse graining}

We often account for tensor networks, which are not defined on a tree hypergraph.
In that case, one can perform graph manipulations, which we call coarse graining schemes, to reformulate the contraction on an equivalent tree hypergraph.

\begin{definition}
    Let $\graph=(\nodes,\edges)$ and $\secgraph=(\node,\secedges)$ be two hypergraphs sharing their node set and $c: \edges \rightarrow \secedges$ a map.
    We say that $c$ is a coarse graining map and $\secgraph$, if for all $\edgein$
    \begin{align*}
        \edge \subset c(\edge) \, .
    \end{align*}
    For any tensor network $\graphtnetwith$ on $\graph$ we further define a coarse grained tensor network $\sectnetofat{\secgraph}{\nodevariables}$ on $\secgraph$ with hypercores to $\secedge\in\secedges$ by
    \begin{align*}
        \sechypercoreofat{\secedge}{\catvariableof{\secedge}}
        = \srcontractionof{\{\edgehypercorewith\wcols\secedge=c(\edge)\}}{\catvariableof{\secedge}} \, .
    \end{align*}
\end{definition}

In the literature such coarse graining procedures appear as Cluster Graphs.

We notice that coarse graining of tensor networks preserves their contractions, since
\begin{align*}
    \graphtnetwith
    &= \srcontractionof{\{\edgehypercorewith\wcols\edgein\}}{\nodevariables} \\
    &= \srcontractionof{\bigcup_{\secedge\in\secedges}\{\edgehypercorewith\wcols c(\edge)=\secedge\}}{\nodevariables} \\
    &= \srcontractionof{\{\srcontractionof{\{\edgehypercorewith\wcols c(\edge)=\secedge\}}{\edgevariables}\wcols\secedge\in\secedges\}}{\nodevariables} \\
    &= \srcontractionof{\{\sechypercoreofat{\secedge}{\catvariableof{\secedge}}\wcols\secedge\in\secedges\}}{\nodevariables} \\
    &= \sectnetofat{\secgraph}{\nodevariables} \, .
\end{align*}

\subsubsection{Variable Elimination}\glsadd{mp:VariableElimination}

Given an order of the nodes the variable elimination procedure provides a coarse graining scheme of an augmented graph.


\begin{algorithm}[hbt!]
    \caption{Variable Elimination}\label{alg:variableElimination}
    \begin{algorithmic}
        \Require Tensor network $\extnetasset$ with total order of the nodes
        \Ensure Coarse graining map $c$ towards a tree hypergraph $(\nodes,\secedges)$

        \iosepline
        \State $\mathcal{C}\algdefsymbol\varnothing$ a set of cores already contracted
        \State $\secgraph=(\varnothing,\varnothing)$ a augmented hypergraph to be constructed
        \State $c:\varnothing\rightarrow\varnothing$ a coarse graining map to be constructed
        \State $\mathcal{M}=\varnothing$ auxiliary edges (messages)

        \For{$\nodein$ following the total order}

            \State {\bf Step 1:} Plan new message $\edgeof{\node}$ given the current ones $\mathcal{M}$
            \begin{itemize}
                \item Choose all edges from $\edges$ containing $\node$, which have not been contracted before and update the contracted ones
                \begin{align*}
                    \edgesof{\node} &\algdefsymbol \{\edge \wcols \edge\in\edges\backslash\mathcal{C} \ncond \node\in\edge \} \\
                    \mathcal{C} & \algdefsymbol \mathcal{C} \cup \edgesof{\node}
                \end{align*}
                \item Choose and remove all message edges from $\mathcal{M}$ containing $\node$
                \begin{align*}
                    \mathcal{M}^{\node} & \algdefsymbol \{\edgeof{\secnode} \wcols \edgeof{\secnode} \in \mathcal{M}\ncond \nodeof{\secnode}\in\edgeof{\secnode} \} \\
                    \mathcal{M} & \algdefsymbol \mathcal{M}\backslash\mathcal{M}^{\node}
                \end{align*}
                \item Build a coarse grained edge and add it to the messages
                \begin{align*}
                    \edgeof{\node} &\algdefsymbol
                    \{\node\} \cup
                    \bigcup_{\edge\in\edgesof{\node}} \{\thirdnodeof{\node} \wcols \thirdnode\in\edge\backslash\{\node\} \} \cup
                    \bigcup_{\edgeof{\secnode}\in\mathcal{M}^{\node}} \{\thirdnodeof{\secnode} \wcols \thirdnodeof{\secnode}\in\edgeof{\secnode}\backslash\{\secnode\} \} \\
                    \mathcal{M} &\algdefsymbol \mathcal{M} \cup \{\edgeof{\node}\}
                \end{align*}
            \end{itemize}

            \State {\bf Step 2}: Build manipulated cores of an augmented network, which are coarse grained to $\edgeof{\node}$
            \begin{itemize}
                \item  Rename the variables at each $\edge\in\edgesof{\node}$, keeping only $\node$
                \begin{align*}
                    \sechypercoreofat{\edge}{\catvariableof{\node},\catvariableof{\secnodeof{\node}}\wcols \secnode\in\edge\backslash\{\node\}} \algdefsymbol
                    \srcontractionof{
                        \hypercoreofat{\edge}{\catvariableof{\edge}}
                        \cup \{\identityat{\catvariableof{\secnode},\catvariableof{\secnodeof{\node}}} \wcols \secnode\in\edge\backslash\{\node\}\}
                    }{
                        \catvariableof{\node},\catvariableof{\secnodeof{\node}}\wcols \secnode\in\edge\backslash\{\node\}
                    }
                \end{align*}
                \item Add delta core to current variable
                \begin{align*}
                    \identityat{\catvariableof{\node},\catvariableof{\nodeof{\secnode}}\wcols\edgeof{\secnode}\in\mathcal{M}^{\node}} \\
                    c(\{\node\}\cup\{\nodeof{\secnode}\wcols\edgeof{\secnode}\in\mathcal{M}^{\node}\}) = \edgeof{\node}
                \end{align*}
                \item Add delta cores to any other variable $\thirdnodeof{\node}\in\edgeof{\node}\{\node\}$
                \begin{align*}
                    \identityat{\catvariableof{\thirdnodeof{\node}},\catvariableof{\thirdnodeof{\secnode}}\wcols\edgeof{\secnode}\in\mathcal{M}^{\node}\ncond\thirdnode\in\secnode}
                    c(\{\thirdnodeof{\node}\}\cup\{\thirdnodeof{\secnode}\wcols\edgeof{\secnode}\in\mathcal{M}^{\node}\ncond\thirdnode\in\secnode\}) = \edgeof{\node}
                \end{align*}
            \end{itemize}
        \EndFor
        \State \Return Augmented tensor network $\sechypercoreof{\secgraph}$ and coarse graining map $c$
    \end{algorithmic}
\end{algorithm}

The variable elimination algorithm is simultaneously defining an augmentation of a tensor network with delta cores and defining a coarse graining mechanism of the augmentation.


\begin{lemma}
    The \algoref{alg:variableElimination} returns a coarse graining map $c$ towards a tree hypergraph.
\end{lemma}
\begin{proof}
    By construction we have $\edge\subset c(\edge)$ for any $\edgein$, thus $c$ is a coarse graining map.

\end{proof}

\begin{figure}
        % Example
    \stantikzpic{
    \drawtextnode{-5}{1}{$a)$}

    \drawsmalltensor{0}{0}{$\hypercoreof{0}$}
    \drawvwire{0}{-1}{$\catvariableof{0}$}
    \drawsmalltensor{0}{-4}{$\hypercoreof{1}$}
    \drawvwire{0}{-5}{$\catvariableof{1}$}
    \drawsmalltensor{0}{-8}{$\hypercoreof{2}$}
    \drawvwire{0}{-9}{$\catvariableof{2}$}

    \drawtextnode{0}{-12}{\colorlabelsize $\vdots$}

    \drawvwire{0}{-13}{$\catvariableof{\catorder-2}$}
    \drawsmalltensor{0}{-16}{$\hypercoreof{\catorder-1}$}


    \drawvariabledot{-4}{-8}
    \draw (-4,-8) to [bend right=-20] (-1,0);
    \draw (-4,-8) to [bend right=-10] (-1,-4);
    \draw (-4,-8) to (-1,-8);
    \draw (-4,-8) to [bend right=20] (-1,-16);

    \begin{scope}[shift={(15,0)}]
        \drawtextnode{-5}{1}{$b)$}

        \drawsmalltensor{0}{0}{$\hypercoreof{0}$}
        \drawvwire{0}{-1}{$\catvariableof{0}$}
        \drawsmalltensor{0}{-4}{$\hypercoreof{1}$}

        \draw (-4,-3) to [bend right=-30] (-1,0);
        \drawsmalltensor{-4}{-4}{$\identity$}
        \drawhwire{-3}{-4}{}
        \drawvwire{-4}{-5}{}

        \drawtextnode{5}{-4}{Cluster to $\catvariableof{0}$}
        \draw[dashed] (-6,-6) -- (8,-6);

        \drawvwire{0}{-5}{}
        \drawsmalltensor{0}{-8}{$\hypercoreof{2}$}
        \drawvwire{0}{-9}{}

        \drawsmalltensor{-4}{-8}{$\identity$}
        \drawhwire{-3}{-8}{}
        \drawvwire{-4}{-9}{}

        \drawtextnode{5}{-8}{Cluster to $\catvariableof{1}$}
        \draw[dashed] (-6,-10) -- (8,-10);

        \drawtextnode{-4}{-12}{\colorlabelsize $\vdots$}
        \drawtextnode{0}{-12}{\colorlabelsize $\vdots$}

        \draw[dashed] (-6,-14) -- (8,-14);
        \drawtextnode{5}{-16}{Cluster to $\catvariableof{\catorder-2}$}

        \drawvwire{-4}{-13}{}
        \drawsmalltensor{-4}{-16}{$\identity$}
        \drawhwire{-3}{-16}{}

        \drawvwire{0}{-13}{}
        \drawsmalltensor{0}{-16}{$\hypercoreof{\catorder-1}$}

        \draw[dashed] (-6,-18) -- (8,-18);
        \drawtextnode{5}{-20}{Cluster to $\decvariable$}

        \drawvwire{-4}{-17}{}
        \drawsmalltensor{-4}{-20}{$\identity$}
    \end{scope}

    \begin{scope}[shift={(30,0)}]
        \drawtextnode{-5}{1}{$c)$}

        \drawtensorblock{-2}{-4}{3}{$\sechypercoreof{\edgeof{0}}$}
        \drawvwire{-4}{-5}{$\decvariableof{\catvariableof{0}}$}
        \drawvwire{0}{-5}{$\catvariableof{0,\catvariableof{1}}$}

        \drawtensorblock{-2}{-8}{3}{$\sechypercoreof{\edgeof{1}}$}
        \drawvwire{-4}{-9}{$\decvariableof{\catvariableof{1}}$}
        \drawvwire{0}{-9}{$\catvariableof{1,\catvariableof{2}}$}

        \drawvwire{-4}{-13}{$\decvariableof{\catvariableof{\catorder-3}}$}
        \drawvwire{0}{-13}{$\catvariableof{\catorder-2,\catvariableof{\catorder-3}}$}
        \drawtensorblock{-2}{-16}{3}{$\sechypercoreof{\edgeof{\catorder-1}}$}

        \drawvwire{-4}{-17}{$\decvariableof{\catvariableof{\catorder-2}}$}
        \drawtensorblock{-4}{-20}{1}{$\sechypercoreof{\decvariable}$} % \edgeof{\decvariable} ?
    \end{scope}

}
    \caption{Variable Elimination on a tensor network in the order $(\catvariableof{0},\ldots,\catvariableof{\catorder-1},\decvariable)$.
    a) Original contraction on a non-tree tensor network, b) Coarse grained Bethe hypergraph and c) Coarse grained contraction.
    }
\end{figure}

\subsubsection{Junction tree algorithm}

Classically the junction tree algorithm is formulated on the extended graph (we define extended graphs here by two nodes adjacent if and only if they are contained in a hypergraph) \cite{lauritzen_local_1988}, where one exploits that a graph has a junction tree if and only if it is triangulated \cite{lauritzen_graphical_1996}.
The junction tree algorithm therefore constructs a triangulation of the graph by adding edges (also called chordal graph \cite{koller_probabilistic_2009}).
The maximal cliques of the triangulated graph form a clique tree, which satisfies the running intersection property \cite{koller_probabilistic_2009}.
Here the edges of the tree can be selected by minimal spanning algorithms.
The coarse graining map to a tree hypergraph is thus the selection of the maximal clique in the triangulated extended graph, which contains the nodes of the hyperedge.
Note that finding a minimal triangulation is $\mathrm{NP}$-hard.

% Tensor network formalism
\todo{We can exploit the running intersection property and add delta tensors copying the variables:
    \begin{itemize}
        \item Assign each variable a cluster $\secedge$ containing the variable
        \item Rename the variables $\nodeof{\edge}$ to each cluster, except those variables assigned to a cluster
        \item Add delta tensors representing the sepsets
    \end{itemize}
    This turns a clique tree into a tree hypergraph.
    The running intersection property ensures, that the augmented tensor network coincides with the original.
}

\begin{figure}
        % Example
    \stantikzpic{
    \drawtextnode{-5}{1}{$a)$}

    \drawsmalltensor{0}{0}{$\hypercoreof{0}$}
    \drawvwire{0}{-1}{$\catvariableof{0}$}
    \drawsmalltensor{0}{-4}{$\hypercoreof{1}$}
    \drawvwire{0}{-5}{$\catvariableof{1}$}
    \drawsmalltensor{0}{-8}{$\hypercoreof{2}$}
    \drawvwire{0}{-9}{$\catvariableof{2}$}

    \drawtextnode{0}{-12}{\colorlabelsize $\vdots$}

    \drawvwire{0}{-13}{$\catvariableof{\catorder-2}$}
    \drawsmalltensor{0}{-16}{$\hypercoreof{\catorder-1}$}


    \drawvariabledot{-4}{-8}
    \draw (-4,-8) to [bend right=-20] (-1,0);
    \draw (-4,-8) to [bend right=-10] (-1,-4);
    \draw (-4,-8) to (-1,-8);
    \draw (-4,-8) to [bend right=20] (-1,-16);

    \begin{scope}[shift={(15,0)}]
        \drawtextnode{-5}{1}{$b)$}

        \draw (0,0) ellipse (3cm and 1cm);
        \drawtextnode{0}{0}{$\decvariable,\catvariableof{0}$}
        \draw (0,-1) -- (0,-1.5);
        \draw[dashed] (-1,-2.5) rectangle (1,-1.5);
        \drawtextnode{0}{-2}{\colorlabelsize $\decvariable,\catvariableof{0}$}
        \draw (0,-2.5) -- (0,-3);
        \draw (0,-4) ellipse (3cm and 1cm);
        \drawtextnode{0}{-4}{$\decvariable,\catvariableof{0},\catvariableof{1}$}
        \draw (0,-5) -- (0,-5.5);
        \draw[dashed] (-1,-6.5) rectangle (1,-5.5);
        \drawtextnode{0}{-6}{\colorlabelsize $\decvariable,\catvariableof{1}$}
        \draw (0,-6.5) -- (0,-7);
        \draw (0,-8) ellipse (3cm and 1cm);
        \drawtextnode{0}{-8}{$\decvariable,\catvariableof{1},\catvariableof{2}$}

        \draw (0,-9) -- (0,-9.5);
        \draw[dashed] (-1,-10.5) rectangle (1,-9.5);
        \drawtextnode{0}{-10}{\colorlabelsize $\decvariable,\catvariableof{2}$}
        \draw (0,-10.5) -- (0,-11);

        \drawtextnode{0}{-11.75}{\colorlabelsize $\vdots$}

        \draw (0,-13) -- (0,-13.5);
        \draw[dashed] (-1,-14.5) rectangle (1,-13.5);
        \drawtextnode{0}{-14}{\colorlabelsize $\decvariable,\catvariableof{\catorder\shortminus2}$}
        \draw (0,-14.5) -- (0,-15);

        \draw (0,-16) ellipse (3cm and 1cm);
        \drawtextnode{0}{-16}{$\decvariable,\catvariableof{\catorder-2}$}

    \end{scope}

}
    \caption{Junction tree on a tensor network in the order $(\catvariableof{0},\ldots,\catvariableof{\catorder-1},\decvariable)$.
    $a)$ Example chain hypergraph, which is already triangulated.
    $b)$ Clique tree to the network.
    }
\end{figure}